% ID: SEL-D0046-OA-E01
% ESTADO: validada. S047.
\textbf{a)} Escribimos $A=I+N$, con $N=\begin{pmatrix}0&a\\0&0\end{pmatrix}$ y $N^2=O$. En la expansión de $(I+N)^n$ sólo sobreviven los términos $I+nN$, pues los demás contienen $N^2$. Así,
\[
A^{2014}=\boxed{\begin{pmatrix}1&2014a\\0&1\end{pmatrix}}.
\]
\textbf{b)} Con $a=2$, $A^3=\begin{pmatrix}1&6\\0&1\end{pmatrix}$, cuya inversa es $\begin{pmatrix}1&-6\\0&1\end{pmatrix}$. Multiplicamos por la izquierda para conservar el orden del producto:
\[
X=(A^3)^{-1}4B
=\begin{pmatrix}1&-6\\0&1\end{pmatrix}\begin{pmatrix}2&0\\3&0\end{pmatrix}
=\boxed{\begin{pmatrix}-16&0\\3&0\end{pmatrix}}.
\]
